\documentclass[../../../include/open-logic-section]{subfiles}

\begin{document}
\olfileid[tr]{sth}{card-arithmetic}{expotough}

\olsection[Bazı Sadeleştirmeler]{Kardinal Üs Alma İşleminde Bazı Sadeleştirmeler}

$\canonord$ bağıntısını tanımlamak biraz uğraştırıcı olsa da sonuç,
kardinal toplama ve çarpma hakkında yararlı bir sonuç olan
\olref[simp]{cardplustimesmax} idi. Bununla birlikte sonlu ötesi üs alma
bu kadar dolaysız biçimde sadeleştirilemez. Nedenini açıklamak için sonlu
durumdan tanıdık bir örüntüyü genişleten bir sonuçla başlıyoruz (ispatı
yüksek bir soyutlama düzeyinde olsa da):

\begin{prop}\ollabel{simplecardexpo}
Her $\cardfont{a},\cardfont{b},\cardfont{c}$ kardinal sayısı için
$\cardexpo{\cardfont{a}}{\cardfont{b} \cardplus \cardfont{c}} =
\cardexpo{\cardfont{a}}{\cardfont{b}} \cardtimes
\cardexpo{\cardfont{a}}{\cardfont{c}}$ ve
$\cardexpo{(\cardexpo{\cardfont{a}}{\cardfont{b}})}{\cardfont{c}} =
\cardexpo{\cardfont{a}}{\cardfont{b} \cardtimes \cardfont{c}}$.
\end{prop}

\begin{proof}
İlk iddia için bir $f\colon
(\cardfont{b}\disjointsum\cardfont{c})\to\cardfont{a}$ fonksiyonu ele
alın. Her $\beta\in\cardfont{b}$ için
$f_\cardfont{b}(\beta)=f(\beta,0)$ ve her
$\gamma\in\cardfont{c}$ için $f_\cardfont{c}(\gamma)=f(\gamma,1)$
tanımlayarak bunu ``ikiye ayırın''. $f\mapsto
(f_\cardfont{b}\times f_\cardfont{c})$ eşlemesi
$\funfromto{\cardfont{b} \disjointsum \cardfont{c}}{\cardfont{a}} \to
(\funfromto{\cardfont{b}}{\cardfont{a}} \times
\funfromto{\cardfont{c}}{\cardfont{a}})$ bir !!a{bijection} olur.

İkinci iddia için bir $f\colon\cardfont{c}\to
(\funfromto{\cardfont{b}}{\cardfont{a}})$ fonksiyonu ele alın; böylece
her $\gamma\in\cardfont{c}$ için bir
$f(\gamma)\colon\cardfont{b}\to\cardfont{a}$ fonksiyonumuz vardır. Her
$\tuple{\beta,\gamma}\in\cardfont{b}\times\cardfont{c}$ için
$f^*(\beta,\gamma)=(f(\gamma))(\beta)$ tanımlayın. $f\mapsto f^*$
eşlemesi
$\funfromto{\cardfont{c}}{(\funfromto{\cardfont{b}}{\cardfont{a}})}
\to \funfromto{\cardfont{b} \cardtimes \cardfont{c}}{\cardfont{a}}$
bir !!a{bijection} olur.
\end{proof}

Sonsuz kardinal sayılarla uğraşırken
$\cardexpo{\cardfont{a}}{\cardfont{b}}$ değerini hesaplamanın kolay bir
yolunu \emph{isterdik}. İşte bu yönde güzel bir adım:

\begin{prop}\ollabel{cardexpo2reduct}
$2\leq\cardfont{a}\leq\cardfont{b}$ ve $\cardfont{b}$ sonsuzsa
$\cardexpo{\cardfont{a}}{\cardfont{b}} =
\cardexpo{2}{\cardfont{b}}$.
\end{prop}

\begin{proof}
\begin{align*}
	\cardexpo{2}{\cardfont{b}} &\leq
	\cardexpo{\cardfont{a}}{\cardfont{b}}\text{, çünkü $2 \leq \cardfont{a}$}\\
	&\leq \cardexpo{(2^\cardfont{a})}{\cardfont{b}}
	\text{, \olref[opps]{lem:SizePowerset2Exp} gereği}\\
	&= \cardexpo{2}{\cardfont{a} \cardtimes \cardfont{b}}
	\text{, \olref{simplecardexpo} gereği} \\
	&= \cardexpo{2}{\cardfont{b}}
	\text{, \olref[simp]{cardplustimesmax} gereği}
\end{align*}
\end{proof}

\olref[card-arithmetic][opps]{lem:SizePowerset2Exp} gereği
$\cardfont{b}<\cardexpo{2}{\cardfont{b}}$ olduğundan bunu daha fazla
sadeleştirebilmeyi gerçekten beklememeliyiz. Ancak bu,
$\cardfont{b}<\cardfont{a}$ olduğunda
$\cardexpo{\cardfont{a}}{\cardfont{b}}$ hakkında ne söyleyeceğimizi
bildirmez. Elbette $\cardfont{b}$ \emph{sonluysa} ne yapacağımızı biliriz.

\begin{prop}
$\cardfont{a}$ sonsuz ve $0<n\in\omega$ ise
$\cardexpo{\cardfont{a}}{n}=\cardfont{a}$.
\end{prop}

\begin{proof}
$\cardexpo{\cardfont{a}}{n} = \cardfont{a} \cardtimes \cardfont{a}
\cardtimes \ldots \cardtimes \cardfont{a} = \cardfont{a}$;
\olref[simp]{cardplustimesmax} gereği.
\end{proof}
\noindent
Ek olarak başka bazı durumlarda
$\cardexpo{\cardfont{a}}{\cardfont{b}}$ büyüklüğünü denetleyebiliriz:

\begin{prop}
$2\leq\cardfont{b}<\cardfont{a}\leq
\cardexpo{2}{\cardfont{b}}$ ve $\cardfont{b}$ sonsuzsa
$\cardexpo{\cardfont{a}}{\cardfont{b}}=\cardexpo{2}{\cardfont{b}}$.
\end{prop}

\begin{proof}
$\cardexpo{2}{\cardfont{b}}\leq
\cardexpo{\cardfont{a}}{\cardfont{b}}
\leq \cardexpo{(\cardexpo{2}{\cardfont{b}})}{\cardfont{b}} =
\cardexpo{2}{\cardfont{b}\cardtimes\cardfont{b}} =
\cardexpo{2}{\cardfont{b}}$; burada
\olref{cardexpo2reduct} içindeki gibi uslamladık.
\end{proof}
\noindent
Fakat bu noktanın ötesinde işler çok daha incelikli hâle gelir.

\end{document}
